% Figure from MATH 451 notes, section 17. Compile with the notes preamble.
\begin{tikzpicture}
\begin{axis}[nfig, width=0.78\textwidth, height=0.42\textwidth,
  xmin=0, xmax=3.1, ymin=0.15, ymax=1.15,
  xtick={1.12,1.4,1.68}, xticklabels={$x_0\!-\!\delta$,$x_0$,$x_0\!+\!\delta$},
  ytick={0.5115,0.7115,0.9115},
  yticklabels={$f(x_0)\!-\!\varepsilon$,$f(x_0)$,$f(x_0)\!+\!\varepsilon$}, xlabel={}]
  \addplot[draw=none, fill=black!9] coordinates {(0,0.5115) (3.1,0.5115) (3.1,0.9115) (0,0.9115)};
  \draw[figR, thick] (axis cs:1.12,0.5115) rectangle (axis cs:1.68,0.9115);
  \addplot[figB, very thick, domain=0:3.05, samples=250] {0.55+0.35*sin(deg(2.2*x))+0.1*x};
  \addplot[only marks, mark=*, mark size=1.6pt, black] coordinates {(1.4,0.7115)};
  \draw[black!50, dotted] (axis cs:1.4,0.15) -- (axis cs:1.4,0.7115);
  \node[font=\scriptsize, figR] at (axis cs:2.42,0.98) {inside the window the graph};
  \node[font=\scriptsize, figR] at (axis cs:2.42,0.90) {exits through the \emph{sides}};
  \draw[figR!70, -stealth, very thin] (axis cs:1.95,0.90) -- (axis cs:1.66,0.80);
\end{axis}
\end{tikzpicture}
